\section{Objectives}
\label{sec:objectives}

The general objective of this thesis is to determine whether Parkinson's
disease can be classified from FMCW radar micro-Doppler gait spectrograms
under an evaluation a clinician could trust. It is pursued through five
specific objectives:

\begin{itemize}
  \item \textbf{O1.} Characterise the dataset and identify the confounds
        capable of producing spurious classification performance, before
        any model is trained.
  \item \textbf{O2.} Establish a confound-controlled classical baseline
        under subject-independent (leave-one-subject-out) validation,
        including deliberately simple reference predictors, such as one
        that knows only how long each recording lasts, that set the score
        any real model must clearly beat.
  \item \textbf{O3.} Develop and evaluate deep models on windowed
        micro-Doppler signatures: a compact convolutional network, a
        pretrained ResNet-18 with a retrained final layer, and a sequence
        model on velocity envelopes, under the identical protocol.
  \item \textbf{O4.} Verify with interpretability methods (Grad-CAM) what
        the trained networks actually attend to, and test the finding with
        an intervention experiment rather than leaving it as a picture.
  \item \textbf{O5.} Report every result against the explicit null floors
        with honest uncertainty, and reconcile the outcome with the
        published literature by reproducing its validation design on this
        same dataset.
\end{itemize}

The conclusions of Chapter~\ref{ch:discussion} return to these objectives one by one.
